\section{Bayesian intervention analysis for branch campaigns}
\label{exec13:causalimpact-reuse}

A bank that launches a campaign in one branch group may want the daily path of
additional applications, rather than a single coefficient pooled over all
branches. Google CausalImpact is a substantive candidate for this question.
Its Bayesian structural time-series approach models how the treated series
would have evolved without the intervention, using its earlier history and,
when available, contemporaneous control series. The proposed reuse decision
is to retain the R implementation as a reference and prospective baseline,
with an isolated R adapter conditional on a compatible environment and
reviewed behavior. No such adapter or fitted package comparison has been
executed here. This disposition preserves a useful alternative without
turning a forecasting model into a general identification engine.
\cite{exec13_ci_reuse,exec13_ci_r}

Here the package name denotes \texttt{google/CausalImpact}. It does not denote
the aGrUM symbolic object discussed in the identification survey. The
selected R source is version 1.4.1 at the commit beginning
\texttt{0f5bf8d97cfb}, committed on 31 March 2026. Its DESCRIPTION retains a
24 September 2025 package date, but the later commit raises the required R
version to 4.5.0. An unchanged package version is consequently insufficient
to reproduce the environment. The repository contains a three-platform
release-R check workflow and substantive input, model and inference tests;
inspection of those files establishes intended checks, not that the checks
passed on the proposed platform.\cite{exec13_ci_r,exec13_ci_semantics}

\subsection{The target and the identifying restriction}

Let the intervention begin after date \(T_0\), let \(H\) be the declared
post-intervention horizon, and write \(Y_t(0)\) for the outcome that the
treated branch group would have produced without the campaign. The realized
path contrast is \(\Delta_t=Y_t-Y_t(0)\). A posterior for this missing path
uses the observed pre-intervention response and the control history
\(X_{1:T_0+H}\). In the state-space formulation of Brodersen and colleagues,
\begin{align}
 Y_t(0)&=Z_t^{\mathsf T}\alpha_t+\varepsilon_t,\\
 \alpha_{t+1}&=T_t\alpha_t+R_t\eta_t.
\end{align}
Here \(\alpha_t\) is a latent state, \(Z_t\) maps it to the scalar response,
\(T_t\) evolves that state, and \(R_t\) maps state innovations into it.
The observation and state innovations have the specified Gaussian laws;
trend, seasonality and regression components determine the matrices and
their parameters. Posterior simulation integrates over the stated parameter,
state and variable-selection uncertainty. The inherited method supplies a
joint predictive law for the entire missing response path, not merely
separate forecasts at each date.\cite{exec13_ci_brodersen2015}

For causal attribution, the control series must remain unaffected by the
intervention, and the pre-period relationship must continue to describe the
untreated post-period response. Neither an accurate pre-period fit nor a
posterior excluding zero proves these restrictions. National advertising,
shared customer acquisition channels, a branch closure or a simultaneous
pricing change can make the proposed donor history inappropriate. A
post-period control may be available for retrospective analysis while being
unavailable when a campaign is first planned. The proposal must distinguish
these two information sets before using retrospective fit to justify a
prospective decision.\cite{exec13_ci_r,exec13_ci_brodersen2015}

An exact construction shows the consequence. Let \(A_t=0\) for eight
pre-period days and \(A_t=1\) for four campaign days. Set
\(x_t^0=10+t\) and suppose the observed daily application count and control
index satisfy
\[
 X_t=x_t^0+15A_t,\qquad
 Y_t=100+2x_t^0+20A_t.
\]
In this first law, the campaign increases applications by 20 per day and
also changes the control index by 15. The pre-period relationship is exactly
\(Y_t=100+2X_t\). Continuing that relationship into the post-period gives
\[
 Y_t-(100+2X_t)=20-2(15)=-10.
\]
The same twelve observed rows arise under a second law in which the
exogenous control path is the observed \(X_t\), the campaign leaves it
unchanged, and the response equation is
\(Y_t=100+2X_t-10A_t\). Under that law the effect is genuinely negative.
The exogenous control path is allowed to shift on the campaign date; no
random timing assumption excludes that possibility in this construction.
Thus the identical observations, including perfect pre-period fit, support
opposite effects under different control-invariance assumptions. This is a
project counterexample to automatic attribution, not a fitted CausalImpact
result or a claim that its algorithm returns the exact linear forecast.
\cite{exec13_cimpact_examples}

More generally, with pre-period slope \(\beta\), donor contamination
\(\gamma A_t\) and response effect \(\tau A_t\), the continued linear
contrast is \(\tau-\beta\gamma\). This decomposition makes a sensitivity
analysis interpretable: an owner can state a plausible range for donor
spillover and inspect when the sign changes. It does not estimate that range
from the same ambiguous history. Similarly, an unmodeled untreated
post-period shift \(\delta\) would enter the contrast as
\(\tau+\delta-\beta\gamma\). A richer state model can represent some
changes, but cannot learn from missing untreated outcomes which part of a
coincident break was caused by the campaign.

For an investment illustration, stipulate a net value of USD 1 per additional
application and a four-day campaign cost of USD 20. The two laws imply
\(4(20)-20=60\) and \(4(-10)-20=-60\) dollars. These are deliberately
specified values, not measured bank economics. Their role is to show why
obtaining a defensible unaffected control or a randomized comparison can
matter more than reducing Monte Carlo error around the wrong counterfactual.
\cite{exec13_cimpact_examples}

\subsection{What the selected implementation actually summarizes}

The R data path takes the response in its first column and controls in later
columns. It checks numeric input, requires more than three total rows and a
pre-period spanning at least three row positions, and rejects missing
covariates. Some out-of-range period endpoints are clipped with warnings.
The automatic path standardizes using pre-period moments and masks the
response after the pre-period before fitting. The supplied-model path instead
requires a separately fitted \texttt{bsts} model with the post-treatment
response missing. These details are part of the computation's information
boundary; an adapter must not silently train on the response it is supposed
to predict counterfactually.\cite{exec13_ci_semantics}

The automatic model includes a local level and optionally seasonality and
regression. The inspected defaults include 1,000 iterations, standardization,
local-level prior scale 0.01, no seasonal component and static regression.
The static regression construction also supplies an expected model size of
three, expected \(R^2=0.8\), and prior degrees of freedom 50. These are
modeling choices to expose and examine, rather than bank-independent
constants that need no justification. The source passes \texttt{seed=1} to
its \texttt{bsts} calls, while predictive observation-noise simulation uses R
random generation. Reproducibility therefore needs the full runtime and
random-state contract; preserving only a displayed seed is insufficient.
The number of iterations is not a convergence certificate.
\cite{exec13_ci_semantics}

The distinction between latent-state means and predictive draws matters.
The implementation averages state contributions for point predictions, but
uses trajectories including simulated observation noise for predictive
quantiles. For draw \(b\), the reported effect transformation is
\[
 \Delta_t^{(b)}=Y_t-\widetilde Y_t^{(b)},\qquad
 C_h^{(b)}=\sum_{t=T_0+1}^{T_0+h}\Delta_t^{(b)}.
\]
Quantiles of \(C_h^{(b)}\) retain dependence between dates. They are not
obtained by adding the marginal lower and upper bounds. The source also
forces cumulative impact to zero before the post-period begins; that
convention must not be misreported as an empirical pre-treatment causal
test.\cite{exec13_ci_semantics,exec13_ci_brodersen2015}

For example, consider four artificial two-day predictive draws:
\[
(80,120),\quad(120,80),\quad(80,120),\quad(120,80).
\]
Each marginal central 95 percent
empirical interval is \([80,120]\), whereas every path sums to 200. The
interval of sums is \([200,200]\); adding marginal endpoints gives
\([160,240]\). If the observed two-day total is 220, the corresponding
draw-wise cumulative effect is always 20. This finite arithmetic checks the
aggregation rule only. It is neither an MCMC experiment nor a coverage result.
\cite{exec13_cimpact_examples}

Relative effects require another distinction. The inspected summary computes
draw-wise relative totals
\[
 R^{(b)}=\frac{\sum_tY_t}{\sum_t\widetilde Y_t^{(b)}}-1
\]
and then summarizes those ratios. A ratio of posterior means can differ
sharply. With actual total 10 and three artificial counterfactual totals
\(-1,1,2\), the ratios are \(-11,9,4\), whose mean is \(2/3\); dividing
10 by the mean counterfactual total and subtracting one gives 14. Zero
denominators are undefined, and small or negative denominators can make a
percentage difficult to interpret. The proposed report retains absolute
effects and explicit denominator diagnostics rather than clipping such
values into a plausible-looking percentage.\cite{exec13_ci_semantics,exec13_cimpact_examples}

The field called \texttt{p} is a one-sided simulation tail-area calculation.
The inspected implementation appends the observed total to the predictive
totals and takes the smaller inclusive tail proportion. For \(B\) predictive
draws its smallest positive value is \(1/(B+1)\); the four-draw construction
above gives \(1/5\), not zero. Its meaning depends on the specified predictive
model and simulation. It is not a design-based randomization probability, a
proof that the causal assumptions hold, or a posterior probability that a
point-null model is true. An interface should preserve the source field and
explain this definition instead of converting every uncertainty measure into
the same generic significance label.\cite{exec13_ci_semantics,exec13_cimpact_examples}

Finally, units determine useful aggregation. Brodersen and colleagues
distinguish a cumulative flow effect from an average stock effect. Daily
additional applications can be summed over the campaign. An extra USD 10,000
of deposit balance observed on each of three dates is USD 10,000 at the last
date; summing those snapshots produces USD 30,000-days under a daily exposure
convention, not USD 30,000 of new deposits or revenue. Funding value would
need a specified balance-time and margin calculation. The adapter request
therefore includes whether the response is a flow, a stock, a rate or an
index, its units and its aggregation rule.
\cite{exec13_ci_brodersen2015,exec13_cimpact_examples}

\subsection{Calendars, failures and a proposed adapter boundary}

The source saves the original time labels and then replaces the internal
time index with sequential row numbers. Its own comment notes that this
precludes automatic filling of missing time points. The missing-date tests
check that omitted labels remain omitted from the output. A Friday and the
following Monday may therefore be neighboring model steps in a sparse input.
For a bank, business-day dynamics and calendar-day accrual are different
choices. Before fitting, the proposed adapter validates a declared calendar,
row uniqueness, time-zone treatment and period boundaries. Filling omitted
rows, carrying forward balances or aggregating transactions would require
an explicit transformation; each can change the estimand or state evolution.
\cite{exec13_ci_semantics}

Constant, entirely missing or insufficiently observed response series cause
the selected model builder to warn and return no fitted model. The wrapper
can still return an object with the CausalImpact class and unavailable
inferences. Checking only the object's class would misclassify that failure
as success. The proposed output boundary instead requires a fitted model,
usable draws and the requested summaries, while retaining warnings and
missing-value reasons. Missing controls, inadmissible periods, failed
inference, unavailable relative summaries and an unsupported scientific
request need distinguishable responses. A treatment-affected donor is a
scientific restriction to assess; it is not equivalent to a numerical error.
\cite{exec13_ci_semantics,exec13_ci_reuse}

For a future R worker, the request should bind the response definition,
units, calendar, pre/post periods, donor identities, supplied causal
restrictions, data provenance, missingness policy, model specification,
priors, iteration and random-state settings, and a resource limit. The
response should retain source/runtime identities, effective periods,
warnings, posterior draw identity or a verifiable draw artifact, point and
cumulative summaries with their uncertainty type, denominator diagnostics
and execution status. Neither the generated prose report nor a generic
``success'' field substitutes for these quantities. Inputs and outputs would
cross a versioned file or process boundary; the existing CLI and MCP service
does not acquire this behavior merely because the contract is specified.
\cite{exec13_ci_reuse}

The first behavioral comparison should run the same frozen fixture directly
in R and through the adapter in the same environment. It should compare
effective periods, row labels, warnings, draw dimensions and numerical
summaries, not only a plotted effect. Separate fixtures should cover a missing
date, missing controls, a constant response, a one-day post-period, a gap
between periods, out-of-range endpoints and a supplied model with incorrectly
retained post-treatment outcomes. The exact constructions above provide
independent checks of aggregation and reporting; they do not replace a
matched native run. A subsequent prospectively designed simulation can
assess recovery under an explicit observation law and retain contaminated
donors and structural breaks as failures. None of these future fitted
comparisons is claimed as completed.\cite{exec13_ci_reuse}

\subsection{Language paths, licenses and the reuse decision}

The R package declares Apache 2.0. Its license requires preserving the
applicable license and attribution notices, marking modifications and
carrying forward any applicable NOTICE material when distributing derivative
work. Its dependency stack is not thereby assigned the same license. The
inspected cutoff CRAN mirrors contain \texttt{bsts} 0.9.11,
\texttt{Boom} 0.9.17 and \texttt{BoomSpikeSlab} 1.2.7. Their descriptions
declare compilation requirements, R constraints and, for the latter two,
LGPL 2.1 alternatives tied to their license files. The \texttt{bsts}
declaration is \texttt{LGPL-2.1 | MIT + file LICENSE}; its companion file
supplies the copyright holder and year. The inspected Boom description also
identifies incorporated third-party code. Any distributed build consequently
needs its own dependency and notice inventory, including the applicable
LGPL source or relinking route. These declarations are source evidence, not
a solved dependency lock or blanket redistribution clearance.
\cite{exec13_ci_r,exec13_ci_dependencies}

The selected R dependency requirement is material: the observed workspace
has R 4.1.2 and none of the seven queried CausalImpact/direct or modeling
packages installed. Replacing the pin with an older release would change
the reviewed candidate. A separate R 4.5-or-later environment with pinned
compatible compiled dependencies is the concrete future route. This
environment cost, data/calendar preparation, donor justification, Monte
Carlo diagnostics, adapter fixtures and maintenance belong in the prototype
budget. No measured runtime or monetary integration estimate is available
from this source review.\cite{exec13_ci_reuse,exec13_ci_dependencies}

For Python, the R README recommends \texttt{google/tfp-causalimpact}.
At the separately selected 5 October 2026 commit beginning
\texttt{2bd0a64f7836}, its exported version is 0.2.0, its distribution name
is \texttt{tfp-causalimpact}, and its import namespace is
\texttt{causalimpact}. It exports \texttt{fit\_causalimpact} and distinct
data, model and inference option types. The project declares Python 3.9 or
later, TensorFlow and TensorFlow Probability among its dependencies, and
Apache 2.0 licensing. The declarations do not freeze compatible backend
versions. Its README describes a lineage involving the R package and Dafiti's
Python implementation and explicitly says it is not an officially supported
Google product. The R README aims for close results while allowing
randomized differences. No interchangeability, backend compatibility or
numerical equivalence has been demonstrated here.
\cite{exec13_ci_python,exec13_ci_r}

The preferred proposal path is therefore to reference the R implementation
and evaluate a narrow native wrapper before considering a port. Directly
adopting its estimation internals into the current library would import a
runtime and output contract that has not been exercised. Porting the R
wrapper and its sampler stack into Python or C++ would add responsibility
for numerical behavior, priors, random generation, failure handling and
dependency terms. The existing TFP implementation is a separate candidate
that should be compared before undertaking such a port. Its backend source,
license closure and native comparison remain deferred. Shared import names
are insufficient package identity.\cite{exec13_ci_reuse,exec13_ci_python}

Alternatives should be selected by the question. A simple seasonal forecast
is a transparent predictive baseline; a donor-weight procedure emphasizes
an explicitly restricted comparison group; a panel design is preferable
when its assignment and population assumptions match the decision; and a
randomized rollout can address causal ambiguity that predictive refinement
leaves intact. The earlier survey describes the respective packages. A
comparison should align the branch group, outcome units, horizon, available
information and uncertainty meaning before ranking losses. A package that
answers a different estimand is not a losing implementation of the same
request.

The research owner must justify the donor and stability restrictions and
declare the target. The adapter owner must supply the compatible environment,
dependency notices, calendar and conversion checks, resource handling and
complete failure outputs. A reviewer must assess native-versus-adapter
behavior and the prospectively specified study before any fitted performance
claim is made. The present decision grants no integration capability.
Reconsider it when a compatible pinned environment and matched native oracle
exist, or when a changed scientific question, dependency revision or target
platform changes the appropriate baseline. Unavailable native execution is
local to this candidate and does not prevent work on other proposal branches.
\cite{exec13_ci_reuse}
\label{exec13:causalimpact-reuse-end}
